\section{Observation model and reporting requirement}
\label{sec:model}

Arms are indexed by arrival order, with a fixed order within a batch. Arm $i$
arrives at a deterministic time $\tau_i$ and remains available. Let
\[
 A_t=\{i:\tau_i\le t\},\qquad K_t=|A_t|<\infty.
\]
We assume $A_1\ne\varnothing$. Arm $i$ has an independent potential reward
stream $(X_{i,n})_{n\ge1}$ of independent, identically distributed observations
in $[0,1]$, with mean $\mu_i$. At round $t$ the policy sees current arrivals,
chooses $I_t\in A_t$ from the observed history and independent randomization,
and observes the next entry of that arm's stream. It then issues its report.
The policy does not know means, future arrivals, or the evaluation endpoint.
Conditioning on a reward-independent exogenous random schedule gives the same
conditional guarantees when that extension is needed.

Write $N_i(t)$ for lifetime sample counts and $\widehat\mu_{i,n}$ for the mean
of the first $n$ potential observations. Define
\[
 B_t=\max_{i\in A_t}\mu_i,\qquad
 \Delta_i(t)=B_t-\mu_i,\qquad
 R_T=\sum_{t=1}^T\Delta_{I_t}(t).
\]
Here $R_T$ is random pseudo-regret; $\E R_T$ is its expectation.
All endpoints $T$ below refer to actual physical pulls.

For $0<\varepsilon<1$, the target is
\[
 G_t^\varepsilon=\{i\in A_t:\Delta_i(t)\le\varepsilon\}.
\]
A report partitions $A_t$ into $C_t,D_t,V_t$: certified good, certified bad,
and unresolved. With $0<\delta<1/2$, let
\[
 E_{\rm lab}=\{\text{for all }t:\ C_t\subseteq G_t^\varepsilon
                                  \subseteq A_t\setminus D_t\}.
\]
A policy is $\delta$-reliable if $\Prob(E_{\rm lab})\ge1-\delta$ in every
allowed reward environment. Positive reports may need revision after an
arrival. A sound negative report remains valid because $B_t$ is nondecreasing.
This statement refers to the model's stationary means and permanent availability.

Reliability alone permits every arm to remain unresolved. To measure progress,
fix a boundary width $0<\eta<1$ and define
\begin{equation}
 d_i^\varepsilon(t)=|\Delta_i(t)-\varepsilon|,\quad
 U_t(\eta)=\sum_{i\in A_t}\ind\{i\in V_t,d_i^\varepsilon(t)>\eta\},\quad
 \mathcal A_T(\eta)=\sum_{t=1}^TU_t(\eta).
 \label{eq:area}
\end{equation}
An arm within the boundary band still requires a correct label whenever a
label is issued; its unresolved time is simply exempt from the service loss.
No assumption removes such arms from the model. The parameters have different
roles: $\varepsilon$ defines membership, $\eta$ defines the precision at which
unresolved time is charged, and a policy's service precision will be denoted
by $\rho$.

The main joint requirement at a fixed $T$ is
\begin{equation}
 \E_\theta R_T\le r,\qquad
 \Prob_\theta\{E_{\rm lab},\mathcal A_T(\eta)\le a\}\ge1-\delta.
 \label{eq:joint_protocol}
\end{equation}
It uses an expected reward cost and a joint high-probability reporting cost.
It is stronger than control of per-arm errors, and differs from false discovery
rate control. In particular, $\mathcal A_T=o(T)$ requires the time-average
\emph{number} of unresolved separated arms to vanish. It does not merely
require a small unresolved fraction among all arrived arms.

For a performance family $\Theta$, write $\mathfrak R_\Theta(a)$ for the
infimum of $\sup_{\theta\in\Theta}\E_\theta R_T$ over policies reliable on the
full bounded model and satisfying the second condition of
\eqref{eq:joint_protocol} for every $\theta\in\Theta$. The value is $+\infty$
if no policy is feasible. The lower bounds remain valid if reliability is
required only on the lower-bound subfamily.

\subsection{Admission and retirement in the reward account}

The distinction between statistical reports and reward eligibility is useful
even before introducing a policy. Let $S_t\subseteq A_t$ be an admitted set and
$F_t\subseteq S_t$ its nonempty reward-active subset, with $I_t\in F_t$.
At every deterministic endpoint, including one inside an algorithmic round,
\begin{align}
 R_T={}&\sum_{t\le T}(\max_{i\in F_t}\mu_i-\mu_{I_t})
 +\sum_{t\le T}(B_t-\max_{i\in S_t}\mu_i)\notag\\
 &+\sum_{t\le T}(\max_{i\in S_t}\mu_i-\max_{i\in F_t}\mu_i).
 \label{eq:three_part}
\end{align}
This follows by adding and subtracting the two intermediate maxima at each
pull. The terms are within-active learning, admission delay, and loss of an
admitted benchmark through retirement. The latter two cannot in general be
identified with arrival increments alone. Our checkpoint policy uses
$F_t=S_t=A_t$ for reward selection. Its reporting exclusions consequently
do not create persistent loss of a reward benchmark.
